\documentclass[11pt,aspectratio=169]{beamer}
\usepackage[utf8]{inputenc}
\usepackage[T1]{fontenc}
\usepackage[spanish]{babel}
\usepackage{graphicx}
\usepackage{booktabs}
\usepackage{listings}
\usepackage{xcolor}
\usepackage{hyperref}

\usetheme{Madrid}
\usecolortheme{default}

\lstset{
    basicstyle=\ttfamily\footnotesize,
    keywordstyle=\color{blue},
    commentstyle=\color{gray},
    stringstyle=\color{red},
    breaklines=true,
    showstringspaces=false,
    frame=single,
    numbers=none,
    language=SQL
}

\title[U-Linker Pitch]{U-Linker: Plataforma de Tutoria P2P \\ con Economia de Tokens}
\subtitle{Pitch Final --- Proyecto de Bases de Datos}
\author{P1 · P2 · P3 · P4 · P5 · P6}
\date{}

\begin{document}

% ======================================================================
\begin{frame}
    \titlepage
\end{frame}

% ======================================================================
% SECCION 1 --- P1: Problema y reglas de negocio (1.5 min)
% ======================================================================
\section{Problema y reglas de negocio}

\begin{frame}{El problema}
    \textbf{Que es U-Linker?}
    \begin{itemize}
        \item Plataforma de tutoria entre pares universitarios
        \item Economia de tokens: el conocimiento se intercambia
        \item Sistema automatizado con integridad en la BD
    \end{itemize}
    \vspace{0.3cm}
    \textbf{Datos reales:}
    \begin{itemize}
        \item 123 usuarios, 20 tutores, 15 materias, 30 servicios
        \item 1051 reuniones, 80 resenas, 15 sanciones
        \item 150+ movimientos de tokens
    \end{itemize}
\end{frame}

\begin{frame}{Reglas de negocio criticas}
    \begin{block}{Integridad garantizada desde la base de datos}
        \begin{enumerate}
            \item \textbf{RN-01:} Solo nota $\geq 4.0$ habilita ser tutor de una materia \\ {\footnotesize (CHECK + trg\_validar\_materia\_tutor)}
            \item \textbf{RN-02:} Una resena por estudiante-tutor-materia \\ {\footnotesize (UNIQUE + CHECK 1-5)}
            \item \textbf{RN-03:} Saldo suficiente de tokens para agendar \\ {\footnotesize (fn\_verificar\_saldo + trg\_control\_saldo)}
            \item \textbf{RN-04:} Cancelacion $>$24h $\rightarrow$ reembolso automatico \\ {\footnotesize (Trg\_reserva\_cancelada)}
            \item \textbf{RN-05:} Usuario baneado no agenda ni ofrece servicios \\ {\footnotesize (fn\_es\_usuario\_baneado + 2 triggers)}
            \item \textbf{RN-06:} Calificacion del tutor se recalcula sola \\ {\footnotesize (Trg\_actualizar\_calificacion\_tutor)}
        \end{enumerate}
    \end{block}
    {\small \textbf{Clave:} Estas reglas viven en la BD, no en el backend ni frontend.}
\end{frame}

% ======================================================================
% SECCION 2 --- P2: Decisiones de diseno (2.0 min)
% ======================================================================
\section{Decisiones de diseno}

\begin{frame}{Del E-R al relacional: 10 tablas}
    \textbf{Entidades:} usuario, estudiante, tutor, materia, servicio, reunion, resena, baneos, movimiento\_token, materia\_aprobada\_tutor

    \vspace{0.2cm}
    \textbf{Relaciones clave:}
    \begin{itemize}
        \item \texttt{materia\_aprobada\_tutor} (M:N con atributo \texttt{nota})
        \item \texttt{reunion}: relacion ternaria estudiante--tutor--servicio
        \item \texttt{materia}: auto-relacion recursiva de prerrequisitos
    \end{itemize}
    \begin{block}{Constraints}
        15 FOREIGN KEYs · 6 CHECK · 5+ UNIQUE · 7 triggers · 6 stored procedures · 5 vistas
    \end{block}
\end{frame}

\begin{frame}{Estrategia de jerarquia: table-per-subtype}
    \textbf{\texttt{usuario} (padre)} $\rightarrow$ \texttt{estudiante} (1:1), \texttt{tutor} (1:1)

    \vspace{0.2cm}
    \textbf{Por que esta estrategia?}
    \begin{itemize}
        \item Un usuario puede ser estudiante, tutor o \textbf{ambos} (rol = `Ambos')
        \item Cada subtipo tiene atributos exclusivos (semestre, calif\_promedio)
        \item Las FK apuntan al subtipo correcto, no al padre generico
    \end{itemize}
\end{frame}

\begin{frame}{Normalizacion y desnormalizacion}
    \begin{block}{Esquema en 3FN}
        No hay dependencias transitivas. Cada atributo depende de la llave completa.
    \end{block}
    \begin{block}{Dos desnormalizaciones justificadas}
        \begin{enumerate}
            \item \textbf{\texttt{usuario.saldo\_tokens}} --- mantenido por SPs/triggers. Evita recalcular desde \texttt{movimiento\_token} (auditoria).
            \item \textbf{\texttt{reunion.tokens\_cobrados}} --- snapshot historico del precio. Si el tutor cambia el precio despues, la reunion conserva el valor original.
        \end{enumerate}
    \end{block}
\end{frame}

\begin{frame}{Auto-relacion recursiva en materia}
    \texttt{materia.id\_prerequisito} $\rightarrow$ FK a \texttt{materia.id\_materia}

    \vspace{0.2cm}
    \textbf{Cadena de ejemplo:} Calculo Diferencial $\rightarrow$ Calculo Integral $\rightarrow$ Estadistica

    \vspace{0.2cm}
    Consultada con CTE recursiva en \texttt{sp\_materias\_con\_prerrequisitos}. ON DELETE SET NULL.
\end{frame}

\begin{frame}{Triggers y SPs: logica en la BD}
    \begin{columns}
        \column{0.5\textwidth}
        \textbf{7 Triggers}
        \begin{enumerate}
            \footnotesize
            \item Trg\_actualizar\_calificacion\_tutor
            \item Trg\_reserva\_cancelada
            \item trg\_validar\_materia\_tutor
            \item trg\_validar\_servicio
            \item trg\_validar\_baneo\_reunion
            \item trg\_validar\_baneo\_servicio
            \item trg\_control\_saldo
        \end{enumerate}
        \column{0.5\textwidth}
        \textbf{6 Stored Procedures}
        \begin{enumerate}
            \footnotesize
            \item pr\_agendar\_tutoria
            \item pr\_asignar\_datos\_usuario
            \item sp\_tutorias\_pendientes
            \item sp\_tutores\_disponibles
            \item sp\_ranking\_tutores
            \item sp\_materias\_con\_prerrequisitos
        \end{enumerate}
    \end{columns}
    \vspace{0.2cm}
    \textbf{Mecanismo @seed\_mode:} variable de sesion que desactiva triggers de validacion durante la carga masiva de datos (seed.sql). Al terminar, los triggers vuelven a estar activos.
\end{frame}

% ======================================================================
% SECCION 3 --- P3: Demo consulta 1 (1.0 min)
% ======================================================================
\section{Demo 1: JOIN multiple}

\begin{frame}[fragile]{Consulta 1 --- Detalle de reuniones finalizadas}
    \textbf{Pregunta:} Quienes participaron en las ultimas reuniones?

    \begin{lstlisting}
SELECT r.fecha, u_est.nombre AS estudiante,
       u_tut.nombre AS tutor, s.nombre AS servicio,
       m.nombre AS materia, r.tokens_cobrados
FROM reunion r
JOIN usuario u_est ON r.id_estudiante = u_est.id_usuario
JOIN usuario u_tut ON r.id_tutor = u_tut.id_usuario
JOIN servicio s ON r.id_servicio = s.id_servicio
JOIN materia m ON s.id_materia = m.id_materia
WHERE r.estado = 'Finalizada'
ORDER BY r.fecha DESC LIMIT 10;
    \end{lstlisting}

    \begin{block}{Punto tecnico}
        5 tablas en JOIN. Misma tabla \texttt{usuario} referenciada 2 veces con alias. Las FK apuntan a los subtipos que heredan la PK de \texttt{usuario}.
    \end{block}
\end{frame}

% ======================================================================
% SECCION 4 --- P4: Demo consulta 2 (1.0 min)
% ======================================================================
\section{Demo 2: Window functions}

\begin{frame}[fragile]{Consulta 2 --- Evolucion del flujo de tokens}
    \textbf{Pregunta:} Como crece la economia de la plataforma?

    \begin{lstlisting}
SELECT DATE(fecha) AS fecha,
       SUM(cantidad) AS balance_diario,
       SUM(SUM(cantidad)) OVER (
           ORDER BY DATE(fecha)
       ) AS acumulado_historico
FROM movimiento_token
GROUP BY DATE(fecha)
ORDER BY fecha ASC;
    \end{lstlisting}

    \begin{block}{Punto tecnico}
        Window function \texttt{SUM() OVER (ORDER BY)} calcula el acumulado sin subconsultas. El GROUP BY colapsa por dia, y la window function opera sobre el resultado agregado.
    \end{block}
\end{frame}

% ======================================================================
% SECCION 5 --- P5: Optimizacion con EXPLAIN (1.5 min)
% ======================================================================
\section{Optimizacion con EXPLAIN}

\begin{frame}[fragile]{EXPLAIN --- Antes del indice}
    \begin{lstlisting}
EXPLAIN SELECT * FROM reunion
WHERE fecha BETWEEN '2025-01-01' AND '2025-03-31';
    \end{lstlisting}
    \begin{itemize}
        \item \texttt{type: ALL} --- Full table scan
        \item \texttt{rows: 1051} --- Todas las filas examinadas
        \item \texttt{Extra: Using where}
    \end{itemize}
\end{frame}

\begin{frame}[fragile]{EXPLAIN --- Despues del indice}
    \begin{lstlisting}
CREATE INDEX idx_reunion_fecha ON reunion(fecha);
EXPLAIN SELECT * FROM reunion
WHERE fecha BETWEEN '2025-01-01' AND '2025-03-31';
    \end{lstlisting}
    \begin{itemize}
        \item \texttt{type: range} --- Acceso por rango
        \item \texttt{key: idx\_reunion\_fecha} --- Usa el indice
        \item \texttt{rows: \textasciitilde200} --- Solo filas del rango
    \end{itemize}
\end{frame}

\begin{frame}{Resultado cuantitativo}
    \begin{table}
        \begin{tabular}{lccc}
            \toprule
            \textbf{Metrica} & \textbf{Antes} & \textbf{Despues} & \textbf{Mejora} \\
            \midrule
            Tipo de acceso & ALL (full scan) & range (indice) & --- \\
            Filas examinadas & $\sim$1051 & $\sim$200 & \textbf{5x menos} \\
            \bottomrule
        \end{tabular}
    \end{table}
    \textbf{Otros indices:} \texttt{usuario(email)} UNIQUE, \texttt{resena(id\_tutor)}, \texttt{movimiento\_token(id\_usuario, fecha)}
\end{frame}

% ======================================================================
% SECCION 6 --- P6: Diferenciador tecnico (1.0 min)
% ======================================================================
\section{Diferenciador tecnico}

\begin{frame}{Por que nuestra solucion es superior}
    \begin{enumerate}
        \item \textbf{Integridad en 3 capas}
        \begin{itemize}
            \item CHECK constraints $\rightarrow$ rechazan datos invalidos
            \item Triggers $\rightarrow$ validan condiciones complejas
            \item SPs $\rightarrow$ operaciones atomicas multi-tabla
        \end{itemize}
        \vspace{0.2cm}
        \item \textbf{Automatizacion total}
        \begin{itemize}
            \item Calificacion y reembolsos: automaticos
            \item Si alguien inserta por consola, la BD igual lo rechaza
        \end{itemize}
        \vspace{0.2cm}
        \item \textbf{Sistema completo}
        \begin{itemize}
            \item 10 tablas, 7 triggers, 6 SPs, 5 vistas, 16 consultas
            \item Backend Flask funcional (24/24 tests OK)
            \item Frontend Bootstrap 5
        \end{itemize}
    \end{enumerate}
\end{frame}

\begin{frame}{Lo que aprendimos}
    \begin{block}{Lecciones}
        \begin{enumerate}
            \item La BD es la \textbf{primera linea de defensa} de la integridad
            \item Triggers y SPs mueven logica al motor = \textbf{mas eficiente y seguro}
            \item Desnormalizar solo con \textbf{justificacion medible}
            \item \texttt{EXPLAIN} no es opcional: \textbf{cada indice tiene metrica}
            \item \texttt{@seed\_mode} permite \textbf{poblar datos historicos} sin disparar validaciones
        \end{enumerate}
    \end{block}
    \vspace{0.5cm}
    \centering
    \Large Gracias! Preguntas?
\end{frame}

\end{document}